\documentclass[11pt]{article}
\usepackage[utf8]{inputenc}
\usepackage[T1]{fontenc}
\usepackage{lmodern}
\usepackage[margin=1in]{geometry}
\usepackage{amsmath,amssymb}
\usepackage{textcomp}
\usepackage[dvipsnames]{xcolor}
\usepackage{enumitem}
\usepackage{parskip}
\usepackage{titlesec}
\usepackage[colorlinks=true,linkcolor=blue,citecolor=blue,urlcolor=blue]{hyperref}
\usepackage[round]{natbib}
\usepackage{microtype}
\emergencystretch=1.5em

\titleformat{\section}{\large\bfseries}{\thesection}{1em}{}
\titleformat{\subsection}{\normalsize\bfseries}{\thesubsection}{1em}{}
\titleformat{\subsubsection}{\normalsize\bfseries\itshape}{}{0em}{}

\title{\textbf{Reputation Conditions on Agreement, Not Fairness:\\ Measuring Exploitation and Collusion in Repeated LLM Bargaining}}
\author{}
\date{}

\newcommand{\ph}{\textit{[placeholder --- fill in after experiments are run]}}

\begin{document}
\maketitle
\vspace{-2em}

\begin{quote}
\small\textit{Assembled from the six confirmed stage outputs of an \texttt{auto-research-discussion} run (seed topic B5.5), post research-taste audit (\textcolor{ForestGreen}{\textbf{PASS}}). Status: \textbf{specification / paper draft --- no experiments have been run.} All numbers below are pre-registered design targets, not results.}
\end{quote}

\begin{abstract}
Autonomous LLM agents are increasingly delegated bilateral transactions --- price negotiation, resource allocation, service contracting --- in which an observable history of past outcomes is often assumed to discipline exploitative behavior without heavier enforcement. We show this assumption rests on an unexamined asymmetry: a reputation channel feeds back whether an episode ended in agreement, never whether the agreed split was fair given the surplus actually available, so a persistently exploitative split and a covertly coordinated one are both self-reinforcing fixed points of exactly the mechanism assumed to eliminate them. We build a repeated bilateral bargaining testbed, instantiated from real human negotiation corpora, in which fixed pairs of LLM agents negotiate across many episodes with the prior-outcome ledger toggled on or off. Exploitation is scored against two independently defensible fairness references; collusion-like dynamics are declared only under a three-conjunct criterion that a benign convergence-to-fairness process cannot also satisfy; and content/order-decoupled history encodings isolate genuine strategic response from context-length artifacts. We test which trajectory regime --- decay, persistence, or collusion-like --- dominates for current LLM bargainers, and what this implies for reputation-only safeguards in autonomous agent commerce.
\end{abstract}

\section{Introduction}

Bilateral bargaining between autonomous software agents is no longer a thought experiment. As large language models are delegated price negotiations, resource splits, and service contracts on behalf of a principal, one lightweight safeguard keeps recurring in how these systems are designed: let both sides see each other's history. The intuition is that a visible track record --- who conceded, who held firm, who walked away --- will discipline bad behavior the way reputation disciplines human counterparties in repeated markets. Framed this way, the history channel is itself a governed persistent-state mechanism --- what an agent observes, retains, and can act on across episodes \citep{alwaysonagents2606} --- and its adequacy as a safeguard is exactly a claim about what that state does and does not encode. This intuition is cheap to implement and easy to justify, which is exactly why it is worth checking rather than assuming.

Reputation governance of LLM agent populations is not itself unexamined --- gossip- and rating-based reputation frameworks can sustain cooperation and isolate exploiters in LLM multi-agent systems \citep[RepuNet;][]{repunet2505}, and market simulations show LLM agents autonomously exploiting weaknesses in rating-based reputation governance \citep{truthmarkettwin2605} --- but in both of those settings the reputation signal carries content beyond the outcome itself: ratings, gossip, advertised-claim audits, and the option to sever ties. The bilateral case that lightweight deployments actually implement is different, and it is this case that remains unexamined: a fixed pair's outcome ledger, as typically implemented, feeds back only whether an episode concluded in agreement --- never whether the agreed split was fair given the surplus that was actually on the table. This is not a minor implementation detail. A bargaining pair that reaches the same lopsided split every time, and a pair that reaches a fair-looking split only by covertly taking turns being the one who profits, both keep agreeing --- and an agreement-only feedback signal cannot tell the two apart from a fair, freely negotiated outcome. Persistent exploitation and tacit collusion are, in this precise sense, both fixed points of the very mechanism assumed to correct them.

The nearest existing work approaches this gap from adjacent but distinct angles. Repeated LLM-agent market simulations such as CompeteAI \citep{competeai} capture cumulative advantage and adaptation across many rounds, but reputation there is implicit in aggregate customer feedback rather than a toggleable bilateral signal, and there is no surplus split to score for exploitation. Work on long-horizon trust, such as the LH-DECEPTION benchmark \citep{lhdeception}, tracks how credibility evolves across a history of interactions, but the trust being modeled sits in an asymmetric supervisor-performer hierarchy, not between two peers who each hold reputational leverage over the other. Social-simulation frameworks like SOTOPIA \citep{sotopia} evaluate rich multi-turn negotiation scenarios, but score each scenario as an independent episode with no persistent memory carried into the next one --- the very quantity this paper needs to track. And the closest theoretical anchor, \citet{calvano2020}'s demonstration that independent Q-learning agents in repeated pricing games reliably discover tacit collusion, establishes that repeated interaction \emph{can} produce coordination without communication for trained policies with a fixed value table. Nor is the prompted case still open as an existence question: \citet{fish2024algorithmic} show that prompted, in-context LLM pricing agents autonomously reach and sustain supracompetitive prices in repeated posted-price oligopoly. What neither line settles is the bilateral question this paper asks: their settings are population-level posted pricing with collusion read off price levels against a competitive benchmark --- no fixed-pair surplus split to score for fairness, no toggleable outcome-history channel as the treatment, and no benign convergence-to-fairness process to rule out, because a posted-price game has no fair reference to converge to.

Turning this gap into a measurable question surfaces its own difficulties. There is no single correct definition of a ``fair'' split once each side's private valuation and outside option are taken into account, so any exploitation score is reference-dependent unless that dependence is made explicit rather than assumed away. Individual pairs of LLM agents behave noisily from one bargaining run to the next, so a trend across repeated episodes has to be separated from ordinary sampling variation before it can be attributed to reputation at all. The two natural signatures of covert coordination --- prices converging across independent pairs, and turn-taking in who gets the better half within a pair --- are each, on their own, indistinguishable from an entirely benign process of everyone converging on the fair price. And because an in-context LLM bargainer's memory is exactly however much transcript fits in its prompt, any behavioral drift across episodes is entangled with context growth in a way that never arises for a trained reinforcement-learning policy with fixed parameters --- a confound with no analogue in the trained-policy collusion literature, and one that prompted-agent pricing studies design around with fixed-size history windows rather than isolate \citep{fish2024algorithmic}.

Our approach turns each of these obstacles into a specific design choice rather than a hedge. Every trajectory claim is required to hold under two independently defensible fairness references at once, so reference-dependence becomes a reported sensitivity check instead of a silent assumption. Behavior is tracked across many independently seeded bargaining pairs with a model that separates pair-level noise from condition-level trend, so noise cannot masquerade as --- or mask --- a real effect. Covert coordination is declared only when price convergence and turn-taking co-occur \emph{while the split stays away from the fair reference} --- a third condition that a benign convergence process cannot also satisfy, which is what makes the distinction checkable rather than asserted. And the history a pair sees is deliberately rewritten into three variants --- the full ledger, a length-matched summary, and an order-scrambled version --- so that behavior driven by context growth can be told apart from behavior driven by what the history actually says and in what order.

We evaluate this design across repeated bargaining scenarios drawn from real human negotiation corpora, comparing conditions in which the outcome history is visible against matched conditions in which it is withheld, alongside scripted non-adaptive opponents and a single-episode protocol that isolates what repeated structure itself contributes.

In summary, we contribute:
\begin{itemize}[leftmargin=1.4em]
\item An identification and formalization of the asymmetry underlying reputation-based governance for agent bargaining: outcome-history feedback conditions on agreement, never on fairness, making persistent exploitation and tacit collusion both stable under the mechanism meant to prevent them.
\item A three-conjunct identifiability criterion that separates covert coordination from benign convergence-to-fairness, together with a direct test of whether it is load-bearing rather than redundant with either signature alone.
\item Content/order-decoupled history encodings that isolate genuine strategic use of reputation from context-length artifacts intrinsic to in-context, rather than trained, bargaining agents.
\item A pre-registered test of which trajectory regime --- decay, persistence, or collusion-like --- dominates for current LLM bargainers under a manipulated reputation treatment, with a design that would, if the dominant regime were persistence or collusion-like, license concrete implications for reputation-only safeguards in autonomous agent-to-agent commerce.
\end{itemize}

\section{Problem Definition}

\noindent\textbf{Formal setup.} Consider a pair of agents $(i,j)$ that plays a sequence of episodes $t = 1,\dots,T$. Each episode $t$ is a bargaining scenario with total disposable surplus $S_t > 0$ and private valuations and outside options known to each agent but not to its counterpart. A negotiation of bounded length yields an outcome $o_t$ --- an agreement flag together with realized surplus shares $s_t^i$ and $s_t^j$ satisfying $s_t^i + s_t^j \leq S_t$. A history channel $H_t = \{o_1,\dots,o_{t-1}\}$ is either exposed to both agents before episode $t$ or withheld from both, held identical in every other respect between the two conditions. The object of study is the long-run behavior of the sequence $\{o_t\}$ relative to a fairness reference computed from each episode's known parameters: does the deviation from that reference decay toward zero, persist at a stable nonzero level, or evolve into a pattern consistent with tacit coordination between the pair?

\noindent\textbf{Why this is hard.} Fixing $S_t$ and the private valuations does not fix a unique fairness reference: two standard candidates --- an even split of $S_t$, and the Nash bargaining solution given each side's outside option --- can disagree on both the magnitude and, in asymmetric cases, the direction of a split's deviation from ``fair,'' so any claim about $\{o_t\}$'s trajectory is only as meaningful as its stated reference. The sequence $\{o_t\}$ is also realized through a stochastic dialogue process, so a genuine condition-level trend must be separated from ordinary pair-to-pair variation before it can be attributed to the history channel $H_t$ at all. Worse, two candidate signatures of coordination --- a shrinking spread of $s_t$ across independently sampled pairs, and an above-chance alternation in which side of a pair receives the larger share --- are each, in isolation, exactly what a benign process of converging on the fair reference would also produce; distinguishing coordination from fairness-convergence requires a condition on $\{o_t\}$ that the benign process cannot also satisfy. Finally, because $H_t$ is delivered as part of an LLM's prompt rather than as input to a fixed-parameter policy, its size necessarily grows with $t$, entangling any measured change in behavior with the growth of the context itself unless the informational content of $H_t$ and its raw length are experimentally decoupled.

\section{Method Overview}

At a conceptual level, the study is a single measurement pipeline with five cooperating parts. Bargaining scenarios and agent pairs enter the pipeline; two independently defensible fairness references and a matched-pair statistical model turn the resulting outcome sequence into a condition-level exploitation trend; a compound coordination test turns the same sequence into a collusion-or-not verdict; and a set of controlled history rewrites turns the whole exercise into a claim about strategic response rather than context-length bookkeeping.

The \textbf{dual-reference exploitation scorer} takes each episode's known surplus and valuations and produces two parallel fairness deviations, reporting only claims that both agree on. The \textbf{matched-condition trend model} takes the resulting exploitation series from many independently seeded pairs under both reputation-on and reputation-off conditions and separates a real condition-level trend from ordinary pair-to-pair noise. The \textbf{joint coordination criterion} takes the same series and declares coordination only when cross-pair convergence and within-pair alternation co-occur while the split remains away from the fair reference --- the third condition being what rules out benign fairness-convergence. The \textbf{decoupled history encodings} take the reputation-on condition and re-render its history three ways --- verbatim, length-matched summary, and order-scrambled --- to isolate whether behavior tracks the history's content and order or merely its growing length. The \textbf{randomized condition design} governs how scenarios, roles, and pairs are assigned, ensuring the final trajectory label describes the condition rather than any single path's early lock-in.

\section{Methodology}

\subsection{Dual-Reference Exploitation Scorer}
\textbf{Challenge.} No single fairness baseline is uniquely correct once private valuations and outside options are asymmetric, so any exploitation score risks being an artifact of an arbitrary reference choice.

\textbf{Insight.} The ambiguity cannot be resolved analytically, but it can be turned into an explicit, checkable sensitivity analysis rather than a hidden modeling decision.

\textbf{Method.} Each episode's realized surplus share is scored against two references computed from that episode's known parameters --- an even split of the total surplus, and the Nash bargaining solution given each side's outside option --- yielding two parallel exploitation readings per episode. A trajectory claim (decay, persistence, or collusion-like) is only reported when both readings agree; episodes or conditions where the two references disagree are reported as not classifiable under this criterion, rather than resolved by an arbitrary tie-break.

\subsection{Matched-Condition Trend Model}
\textbf{Challenge.} A single bargaining pair's behavior is too noisy to tell a reputation-driven trend apart from ordinary sampling variation in LLM dialogue outcomes.

\textbf{Insight.} Modeling pair-level variation explicitly, rather than pooling or eyeballing raw trajectories, is what lets a condition-level claim survive individual pairs behaving unpredictably.

\textbf{Method.} Many independently seeded pairs are run per condition inside one fixed harness, identical in scenario sampling, turn budget, and prompt scaffolding across reputation-on and reputation-off arms except for the history content itself. A mixed-effects model regresses the exploitation reading on episode index interacted with condition, with random slopes and intercepts per pair, yielding two separately falsifiable estimands from a single fit: the within-condition trend over episodes, and the on-versus-off contrast in that trend. This matched-pair design follows the variance-reduction logic demonstrated for learning-based multi-agent economic simulators by \citet{pairedseeds2512} --- sharing seed-level structure across the reputation-on/off arms, rather than drawing independent seeds for each, is what lets the power analysis below target a standardized slope contrast rather than a cruder mean-difference test.

\subsection{Joint Coordination Criterion}
\textbf{Challenge.} The two natural signatures of covert coordination --- cross-pair price convergence and within-pair alternation of who gets the better share --- are each individually indistinguishable from a benign process of converging on the fair outcome.

\textbf{Insight.} A benign convergence process moves the split \emph{toward} the fair reference; genuine coordination does not have to. Requiring evidence of both signatures \emph{plus} a sustained deviation from the fair reference is therefore exactly the condition a benign process cannot also satisfy.

\textbf{Method.} A condition is declared collusion-like only if (i) the spread of realized shares across independently sampled pairs is shrinking with episode index, (ii) the rate at which the pair's better-off side alternates across episodes exceeds chance, and (iii) the mean absolute deviation from the fair reference remains above a pre-specified tolerance through the final episodes. All three conditions must hold jointly, under both fairness references from the scorer above.

\subsection{Decoupled History Encodings}
\textbf{Challenge.} Because an in-context LLM bargainer's memory is exactly however much transcript is in its prompt, any behavioral change across episodes is entangled with the prompt's growing length --- a confound with no analogue for a fixed-parameter trained policy.

\textbf{Insight.} Content and length grow together in the raw history-on condition, so the only way to separate their effects is to build variants that hold one fixed while varying the other.

\textbf{Method.} Within the reputation-on condition, the same episode sequence is additionally run under a fixed-token-budget summary of the ledger (holding length constant across episodes) and an order-scrambled version of the full ledger (holding length and content constant while destroying temporal order). Comparing trend magnitude across all three variants attributes any observed effect to content, order, or mere length growth. This decomposition mirrors two independent findings on how memory structure, not raw channel capacity, governs repeated-interaction behavior in LLM agents: bounded-memory testbeds that ablate typed memory layers by mutability and retrieval role to isolate which component of history drives behavior \citep{agenticsts2607}, and evidence that memory \emph{architecture} rather than signal volume determines whether repeatedly-interacting LLM agents converge on stable conventions in an unrelated language-emergence setting \citep{memoryarchlangemergence2607} --- both support treating history encoding as a manipulable variable in its own right rather than a fixed implementation detail.

\subsection{Randomized Condition Design}
\textbf{Challenge.} Bargaining is path-dependent --- an early outcome can lock a pair into a regime it never revisits --- so a trajectory label computed from an unrandomized or small sample risks describing one path's history rather than the condition's typical behavior.

\textbf{Insight.} Randomizing the initial conditions that a single path would otherwise lock in converts path-dependence from a confound into noise that averages out across a large enough ensemble.

\textbf{Method.} First-mover role, buyer/seller assignment, and scenario draw are randomized independently across pairs within each condition, and the trajectory classification is computed only at the condition level over a pre-specified minimum number of pairs --- never certified from any single pair's path.

\section{Experiments}

\subsection{Experimental Setup}

\textbf{Datasets.} CraigslistBargain \citep{he2018} --- real single-issue price-negotiation scenarios; used for the pilot stage because a single-issue setting gives the Nash bargaining reference a closed form. Deal-or-No-Deal \citep{lewis2017} --- multi-issue division scenarios with private valuations; used to test generalization beyond single-issue price bargaining. GLEE \citep{shapira2024glee} --- a parameterized economic-game benchmark; used specifically to test whether reference agreement (RQ1 below) depends on controlled valuation asymmetry, rather than only being observed post hoc across the two fixed corpora above. A fourth, human-subjects source --- the behavioral-lab bargaining corpus of \citet{heddaya2023}, which contrasts alternating written offers against unstructured natural-language bargaining over an identical house-price scenario --- is not used to instantiate episodes (it has no repeated-interaction or reputation-channel structure to manipulate) but is the closest empirical anchor for what a single-episode human bargaining distribution looks like, informing the plausibility check on scenario realism in Stage 1 piloting.

\textbf{Baselines} (sanity-check comparators, not competitors to beat --- this is a measurement study, and no research question here is framed as ``our procedure beats baseline X''). A reputation-off arm, identical to the reputation-on arm in every respect except an empty history ledger, is the necessary control for every on-condition claim. Scripted Boulware and tit-for-tat opponents --- non-adaptive and memoryless beyond one episode --- check that any observed trend is not simply generic concession drift against any opponent. A single-episode protocol, adapted from SOTOPIA's \citep{sotopia} episodic evaluation structure (borrowing its protocol, not its reported results), isolates what repeated interaction itself adds relative to one-shot bargaining.

\textbf{Metrics.} The dual-referenced exploitation reading and its mixed-effects trend slope; cross-pair variance of realized shares and its trend; within-pair alternation rate against the chance rate; and the resulting condition-level trajectory classification (decay / persistence / collusion-like / unclassified).

\textbf{Research questions:}
\begin{enumerate}[leftmargin=1.6em]
\item RQ1 --- Do the two fairness references agree on trajectory classification?
\item RQ2 --- Does exploitation decay or persist under reputation visibility?
\item RQ3 --- Does reputation visibility causally shift the exploitation trend relative to no history?
\item RQ4 --- Is the sustained-unfairness conjunct load-bearing for the coordination criterion?
\item RQ5 --- Which coordination signature emerges first, price convergence or turn-taking alternation?
\item RQ6 --- Does the observed trend reflect the history's content and order, or context length alone?
\item RQ7 --- Is the trajectory classification stable across independent random sub-samples of pairs?
\end{enumerate}

\subsubsection{RQ1: Do the two fairness references agree on trajectory classification?}
\textbf{Setup.} For every tested (backbone $\times$ domain) condition cell, compute the trajectory classification twice --- once under the even-split reference, once under the Nash-bargaining-solution reference --- and compare. Sample size: sized to detect disagreement concentrated in the highest-asymmetry quartile of scenarios versus uniform disagreement.

\textbf{Expected outcome.} The two references agree on classification in at least 5/6 of tested cells; where they disagree, disagreement concentrates in cells with the most asymmetric private valuations rather than being spread uniformly.

\textbf{Falsification condition.} Refuted if agreement falls below 5/6 of cells, or if disagreement is spread uniformly across asymmetry levels rather than concentrated in the highest-asymmetry quartile --- either would mean reference choice is a source of arbitrary noise rather than a structured, reportable risk.

\textbf{Hypothesis (plain language).} Reference ambiguity is a real but bounded and structured risk, not a source of arbitrary noise across the whole design.

\textbf{Metrics.} Cross-reference classification agreement rate; asymmetry-conditioned disagreement rate.

\textbf{Result.} \ph

\subsubsection{RQ2: Does exploitation decay or persist under reputation visibility?}
\textbf{Setup.} Pilot stage: 2 backbones $\times$ 2 conditions (reputation on/off) $\times$ 50 pairs $\times$ 20 episodes on CraigslistBargain. Scale-up: $\geq$4 backbones (including cross-model pairs) $\times$ 3 domains $\times$ 100 pairs $\times$ 40 episodes. Mixed-effects model: exploitation $\sim$ episode $\times$ condition $+ (1 + \text{episode} \mid \text{pair})$, likelihood-ratio test on the episode term, $\alpha = 0.05$, Holm-corrected across backbones and the two fairness references.

\textbf{Sample-size rationale.} 50 pairs are sized to detect a standardized slope contrast of $d \approx 0.4$ at $\approx$80\% power; an observed contrast below that in the pilot triggers a pre-committed enlargement to 150 pairs/cell in the scale-up, not a null claim.

\textbf{Expected outcome.} The modal classification under reputation-on is persistence --- a slope not significantly different from zero with final-quintile mean deviation above the pre-registered tolerance ($\varepsilon = 0.10$) --- rather than decay, in the majority of tested cells.

\textbf{Falsification condition.} Refuted toward decay if the majority of cells show a significantly negative slope with final-quintile mean deviation inside the tolerance band under both fairness references; refuted toward collusion if the joint coordination criterion (Methodology) fires instead.

\textbf{Hypothesis (plain language).} Visible outcome history alone does not reliably erase a stable exploitative split within the tested episode horizon.

\textbf{Metrics.} Per-cell mixed-effects slope and 95\% CI; final-quintile mean absolute exploitation deviation; per-cell classification.

\textbf{Result.} \ph

\subsubsection{RQ3: Does reputation visibility causally shift the exploitation trend relative to no history?}
\textbf{Setup.} Same mixed-effects specification as RQ2, reading the episode $\times$ condition (on/off) interaction term rather than the marginal on-condition slope; sign consistency checked across backbones.

\textbf{Expected outcome.} The interaction term is directionally corrective (more negative slope under reputation-on than off, standardized $|d| \geq 0.2$, Holm-corrected $p < 0.05$) but small enough that on-condition cells do not reclassify from persistence to decay --- reputation visibility helps, but under-corrects.

\textbf{Falsification condition.} Refuted (reputation is inert or harmful) if the interaction term is non-significant or positive; refuted in the opposite direction (reputation fully self-corrects) if on-condition cells reclassify to decay while matched off-condition cells do not.

\textbf{Hypothesis (plain language).} Reputation visibility is not inert, but it is not sufficient on its own to fully discipline exploitation either.

\textbf{Metrics.} Interaction-term estimate and CI, pooled and per-backbone; proportion of on-condition cells reclassified relative to matched off-condition cells.

\textbf{Result.} \ph

\subsubsection{RQ4: Is the sustained-unfairness conjunct load-bearing for the coordination criterion?}
\textbf{Setup.} Recompute the coordination verdict with the sustained-unfairness conjunct removed (i.e., using only cross-pair convergence and within-pair alternation), and compare against the full three-conjunct verdict across all tested cells.

\textbf{Expected outcome.} At least one condition that the full criterion classifies as decay or unclassified is reclassified as collusion-like when the third conjunct is dropped, demonstrating that the conjunct changes the verdict rather than being redundant with the other two.

\textbf{Falsification condition.} Refuted if every cell's classification is identical with and without the third conjunct --- that would mean sustained unfairness is redundant with the other two signatures rather than load-bearing.

\textbf{Hypothesis (plain language).} The distinction between ``converged to fair'' and ``colluded'' depends on checking that the split stays unfair, not merely that it stabilizes.

\textbf{Metrics.} Count and identity of cells whose classification changes when the third conjunct is dropped.

\textbf{Result.} \ph

\subsubsection{RQ5: Which coordination signature emerges first, price convergence or turn-taking alternation?}
\textbf{Setup.} Conditional on at least 3 non-decay cells existing (otherwise this question is reported as not evaluable, per protocol --- no substitute analysis). Among non-decay cells where at least one coordination signature reaches significance, compare the episode index at which cross-pair convergence first reaches significance against the episode index at which within-pair alternation first does, via paired comparison across qualifying cells.

\textbf{Expected outcome.} Within-pair alternation reaches significance at least 5 episodes earlier than cross-pair convergence in at least 70\% of qualifying cells, because alternation requires only local two-agent reciprocity while convergence requires the same dynamic to replicate independently across separately seeded pairs.

\textbf{Falsification condition.} Refuted if within-pair alternation fails to reach significance at least 5 episodes before cross-pair convergence in at least 70\% of qualifying cells.

\textbf{Hypothesis (plain language).} If coordination emerges, it is visible at the level of a single pair before it is visible at the level of the population.

\textbf{Metrics.} Onset-episode index per signature per cell; proportion of qualifying cells where alternation precedes convergence.

\textbf{Result.} \ph

\subsubsection{RQ6: Does the observed trend reflect the history's content and order, or context length alone?}
\textbf{Setup.} Within reputation-on, compare trend magnitude under the full ledger, the fixed-length summary, and the order-scrambled ledger, via a three-arm mixed-effects model with encoding as a fixed effect and bootstrap confidence intervals on the magnitude ratios.

\textbf{Expected outcome.} The fixed-length-summary arm retains at least half the full-ledger trend magnitude (ruling out pure context-growth artifact), while the order-scrambled arm retains less than half (showing temporal order, not just content presence, matters).

\textbf{Falsification condition.} Refuted (context-length artifact) if the summary arm's retained magnitude falls below half; refuted (order does not matter) if the scrambled arm retains half or more of the full-ledger magnitude.

\textbf{Hypothesis (plain language).} Whatever behavioral change is observed is a response to what the history says and in what order, not simply to how much text is in the prompt.

\textbf{Metrics.} Slope-magnitude ratios (summary/full, scrambled/full) with bootstrap confidence intervals.

\textbf{Result.} \ph

\subsubsection{RQ7: Is the trajectory classification stable across independent random sub-samples of pairs?}
\textbf{Setup.} Split each condition's pairs into two disjoint random halves and compute the trajectory classification independently on each half; compare agreement across (backbone $\times$ domain) cells.

\textbf{Expected outcome.} The two half-sample classifications agree in at least 80\% of tested cells.

\textbf{Falsification condition.} Refuted if agreement falls below 80\% of cells --- that would mean the trajectory label is sensitive to which pairs happen to be sampled rather than describing the condition itself.

\textbf{Hypothesis (plain language).} The reported trajectory label describes the condition as a whole, not an artifact of which specific pairs happened to be sampled.

\textbf{Metrics.} Split-half classification agreement rate.

\textbf{Result.} \ph

\vspace{0.5em}
\noindent\textbf{Further evaluation.} With more time, a natural extension is to test partner-choice and reputation-visible \emph{markets} (many-to-many, not fixed pairs) --- explicitly out of scope here because it changes the unit of analysis from a dyad's trajectory to a population's, and would need its own identifiability design rather than an extension of the joint coordination criterion above.

\subsection{Threats to Validity}

\textbf{Construct validity.} The exploitation and coordination constructs are only as sound as their extraction and reference-computation instruments. The offer/agreement parser (rule-based with an LLM fallback) is validated against $\geq$2 independent human raters with a Cohen's $\kappa \geq 0.7$ acceptance threshold and an explicit confusion-structure inspection (adjacent-value misreads vs.\ systematic misparse of one offer type), audited on 50 transcripts in the pilot and 300 pre-scale-up; a threshold miss blocks scale-up rather than being absorbed into noisier metrics. A separate construct threat is reference-selection itself: since no single fairness baseline is uniquely correct under asymmetric valuations (Methodology, Dual-Reference Exploitation Scorer), every trajectory claim is required to hold under both the even-split and Nash-bargaining references, with disagreement reported as not classifiable rather than resolved by fiat (RQ1).

\textbf{Internal validity.} Two confounds threaten a causal reading of the reputation-on/off contrast. First, scenario parameters must yield well-defined surplus and outside options for both references to exist at all; scenarios failing static validation are excluded before any API spend, with the exclusion rate reported rather than silently dropped. Second, and more consequentially, the empty ledger in the OFF condition and the fixed-length summary encoding in the history-encoding control must hold prompt format and token budget constant relative to the full-ledger ON condition, or any observed effect is confounded with prompt shape rather than attributable to reputation content; this is checked by an automated per-episode token-count equality assertion in the harness (summary budget fixed at the episode-1 ledger length $\pm$5\%), and a violation is pre-committed to downgrade the causal reading of RQ3 and RQ6 specifically --- the two hypotheses that assertion is designed to protect --- rather than blocking the whole run.

\textbf{External validity.} The study is scoped to sandboxable, two-party bargaining with fully known (to the instrument, not the agents) valuations across three corpora (CraigslistBargain, Deal-or-No-Deal, GLEE); it does not cover population-level markets with partner choice, non-sandboxable live deployments, or bargaining domains without a computable closed-form fairness reference. Findings are further scoped to the backbones and domains actually tested (Stage 1: 2 backbones $\times$ 1 domain; Stage 2: $\geq$4 backbones $\times$ 3 domains) and are not asserted to generalize beyond that grid without further testing.

\textbf{Conclusion validity.} A single pair's trajectory cannot distinguish a reputation-driven trend from ordinary LLM sampling variability, so every trend claim is licensed only through the matched-condition mixed-effects model (many independently seeded pairs, pair-level random slopes) rather than eyeballed from individual trajectories, with a pre-registered power target (50 pairs detect a standardized slope contrast $d \approx 0.4$ at $\approx$80\% power) and a pre-committed enlargement rule --- rather than a null claim --- for underpowered contrasts. A distinct non-degeneracy threat (pairs that always fail to agree, or instantly accept the list price) would starve the exploitation series of usable points regardless of sample size; this is gated before any hypothesis is scored (walking-skeleton gate, $\geq$80\% of episodes ending in non-degenerate agreement), and its failure reopens the skeleton rather than being treated as a substantive null result. Bargaining is also path-dependent --- an early outcome can lock a pair into a regime it never revisits --- so the trajectory label is certified only at the condition level over a pre-specified minimum number of randomized pairs (RQ7's split-half reliability check), never from any single pair's path.

\section{Related Work}

The closest prior work to this paper comes from two directions, and it is worth stating the deltas plainly before surveying the wider landscape. On the dynamics side, \citet{fish2024algorithmic} already show that prompted LLM agents tacitly collude in repeated posted-price oligopoly, and CompeteAI \citep{competeai} already studies repeated LLM-agent interaction with feedback that shapes future behavior. On the reputation side, RepuNet \citep{repunet2505} already shows a content-bearing reputation mechanism sustaining cooperation in LLM agent populations, and the LH-DECEPTION benchmark \citep{lhdeception} already tracks how credibility evolves across a long history. None of the four toggles an agreement-only observable-history channel as a controlled bilateral treatment with a surplus split to score against a fairness reference, which is precisely the identifiability gap this paper is built to close.

\textbf{Repeated-interaction LLM-agent simulations.} CompeteAI \citep{competeai} situates competing LLM agents in a persistent, resource-constrained market and studies cumulative advantage and adaptation over many rounds; SOTOPIA \citep{sotopia} scores rich multi-turn social scenarios, including negotiation, along multiple behavioral dimensions using automated judges. Both establish that repeated or richly-scored LLM-agent interaction produces measurable behavioral dynamics, but neither isolates an observable-history channel as a toggleable treatment with a bilateral surplus split to score --- CompeteAI's reputation is implicit in aggregate market feedback, and SOTOPIA scores each scenario as an independent episode with no persistent cross-episode memory for a fixed pair. A broader survey of persistent memory and state management in LLM agents \citep{alwaysonagents2606} frames exactly this kind of cross-episode channel as governed state with its own lifecycle (observe/write/retrieve/forget/rollback) rather than an incidental logging detail, which is the framing this paper adopts for the reputation ledger specifically.

\textbf{Trust, deception, and strategic behavior under partial information.} The LH-DECEPTION benchmark \citep{lhdeception} models an evolving trust state in a supervisor-performer hierarchy and audits for accumulated deception; AmongUs \citep{amongus} elicits emergent strategic misbehavior in a discrete social-deduction game, scored by Elo rather than by a continuous outcome; StrategicPersuasion \citep{strategicpersuasion} formalizes one-shot sender-receiver persuasion with explicit utilities. Each of these studies asymmetric or single-encounter strategic behavior, but none covers symmetric peer-to-peer bargaining with reputational leverage held by both sides across a genuinely repeated interaction --- LH-DECEPTION's hierarchy has no peer symmetry, AmongUs's win/lose structure cannot express a graded exploitation of a continuous split, and StrategicPersuasion's one-shot games admit no history channel to manipulate in the first place. A related recent finding that LLM agents' public communication diverges from their private, off-the-record stance under relational and audience pressure across repeated multi-agent debate turns \citep{nooneswatchingdebates2607} is thematically adjacent --- both papers study strategic behavior that changes under observation --- but its setting is debate persuasion rather than a bilateral surplus split with a reputation channel to toggle. In the concurrent repeated-game line, When Agents Lie \citep{whenagentslie2607} separates private intent from public announcement in repeated $n$-player games and finds that deviations from announcements are predominantly premeditated --- repeated-interaction exploitation dynamics, but again with no surplus split, no fairness reference, and no reputation treatment to manipulate.

\textbf{Repeated-game collusion and negotiation-dialogue foundations.} \citet{calvano2020} show that independent Q-learning agents in a repeated pricing game reliably discover supra-competitive, tacitly collusive prices without communication --- the standard demonstration that repeated interaction alone can produce coordination. \citet{fish2024algorithmic} extend this demonstration to prompted LLM agents: GPT-4-based pricing agents autonomously sustain supracompetitive prices in the same repeated-oligopoly structure, with seemingly innocuous prompt phrasing modulating the degree of collusion --- but the setting remains population-level posted pricing, collusion is read off price levels against a competitive benchmark rather than through an identifiability criterion, and there is neither a bilateral surplus split to score for exploitation nor a reputation channel to toggle. SocialJax \citep{socialjax} provides a vectorized suite of sequential social-dilemma environments with trained MARL baselines exhibiting related exploitation/cooperation dynamics. A conceptual taxonomy mapping human anti-collusion mechanisms onto multi-agent AI systems \citep{anticollusionmapping2601} directly engages the same algorithmic-collusion concern raised by \citet{calvano2020}, but is prescriptive rather than measurement-based --- it proposes governance mechanisms without itself running a repeated-bargaining instrument to detect when collusion-like dynamics occur. \citet{lewis2017} established the multi-issue negotiation-dialogue task this paper's scale-up domain builds on, training self-play agents on independent dialogues with no cross-dialogue memory, and \citet{heddaya2023} later collected a human-subjects bargaining corpus contrasting alternating-offer and free-form natural-language negotiation over an identical scenario --- but, like Lewis et al.'s dialogues, as single episodes with no repeated-interaction or reputation channel to manipulate. What remains unestablished across all of these is not whether prompted LLM agents can develop repeated-game coordination --- Fish et al.\ settle that for posted pricing --- but whether an observable outcome-history channel corrects, entrenches, or coordinates exploitation of a bilateral surplus split in fixed-pair bargaining, and whether any such dynamic reflects the history's content and order rather than the growth of the transcript itself --- the gap this paper's toggled ledger and encoding controls are specifically built to close.

\textbf{Training environments for strategic language agents.} ARIA \citep{aria} and MARSHAL \citep{marshal} both treat negotiation- and bargaining-adjacent games as training environments --- the former for lower-variance credit assignment via intention-space reward aggregation, the latter for multi-turn self-play policy learning via role-separated advantage estimation. Both measure learning performance rather than the behavioral trajectory of an already-capable agent under a reputation manipulation, and neither includes a reputation-on/off control, since neither is designed to ask whether visibility itself changes behavior.

\section{Conclusion}

This work reframes a governance assumption --- that visible outcome history disciplines bilateral exploitation --- as an empirical question with a precise identifiability structure, rather than treating reputation as self-evidently corrective. The central methodological wager is that decay, persistence, and covert coordination are only distinguishable once fairness is scored against more than one defensible reference and coordination is required to survive a compound, benign-convergence-resistant test; weaker designs risk mistaking either for the other. Whichever regime dominates for current LLM bargainers, the finding is decision-relevant: persistence would mean reputation visibility under-corrects and needs a companion mechanism; collusion-like dynamics would mean visibility itself can be actively counterproductive. Two limitations bound this claim from the outset and are natural next steps rather than afterthoughts: the study is scoped to sandboxable, two-party bargaining with computable valuations, leaving population-level markets with partner choice and non-sandboxable live deployments for follow-up; and the reference-agreement check (RQ1) may itself surface disagreement in high-asymmetry scenarios that this design reports rather than resolves, pointing toward richer fairness formalisms as the next methodological target.

\bibliographystyle{plainnat}
\bibliography{refs}

\vspace{1em}
\hrule
\vspace{1em}

\section*{Pre-handoff checklist}
\begin{itemize}[leftmargin=1.4em]
\item \textbf{Organization/navigation} --- narrative flows problem $\to$ method $\to$ experiments $\to$ related work without forward references to undefined terms.
\item \textbf{Terminology discipline} --- ``fairness reference,'' ``reputation channel,'' ``trajectory classification,'' ``coordination criterion'' each defined at first use in Problem Definition, used consistently thereafter; no internal stage IDs (\texttt{sp\_*}, \texttt{m\_*}, \texttt{rq\_*}) leaked into prose --- checked directly against the discussion pipeline log.
\item \textbf{Notation with worked example} --- the formal setup section defines $(i, j, t, S_t, o_t, H_t)$ and the Problem Definition/Methodology sections consistently reuse this notation; no separate worked numeric example is included since no numbers exist yet (would require fabricating results).
\item \textbf{RQ framing matches outcomes} --- all 7 RQs are graded/quantitative (rates, slopes, agreement proportions), each answered in its own subsection with an explicit falsifiable expected outcome, not a binary yes/no.
\item \textbf{No typesetting artifacts} --- only explicit placeholder markers in Results, per the drafting convention this is a pre-registered spec, not a report.
\item \textbf{Artifact plan} --- the testbed harness and the self-generated transcript corpus are the natural release artifacts; not yet built, so not claimed as available.
\end{itemize}

\section*{Hand-off}
This is a draft, not a finished paper --- experiments have not run. Natural next steps: prose polish and venue-specific LaTeX formatting once a target venue is chosen, or revising an earlier stage if piloting surfaces a problem (e.g., if the reference-agreement rate in RQ1 is much lower than expected).

\end{document}
